% Do not change document class, margins, fonts, etc.
\documentclass[a4paper,oneside,bibliography=totoc]{scrbook}

% some useful packages (add more as needed)
\usepackage[utf8]{inputenc}
\usepackage{graphicx}
\usepackage{latexsym}
\usepackage{amsmath}
\usepackage{amssymb}
\usepackage{tabularx}
\usepackage{booktabs}
\usepackage{algorithm} % you can modify the algorithm style to your liking
\counterwithin{algorithm}{chapter} % so that algorithms have chapter numbers as well
\usepackage{algorithmic}
\usepackage{csquotes}
\renewcommand{\algorithmiccomment}[1]{\hfill\textit{// #1}}
\usepackage[dvipsnames]{xcolor}
\usepackage[colorlinks,citecolor=Green]{hyperref} % you may change/remove the colors
\usepackage{lipsum} % you do not need this

% chicago citation style
\usepackage{natbib}
\bibliographystyle{chicagoa}
\setcitestyle{authoryear,round,semicolon,aysep={},yysep={,}} \let\cite\citep

% example enviroments (add more as needed)
\newtheorem{definition}{Definition} \newtheorem{proposition}{Proposition}

\begin{document}

\frontmatter \subject{Master Thesis} % change to appropriate type
\title{NL2Vis: Multi-agent System for Natural Language to Data Visualization}
\author{Jacob Chen\\
 (matriculation number 1980440)} \date{January 25, 2026}
\publishers{{\small Submitted to}\\
 Social Data Science Group\\
 Prof.\ Dr.\ Marc Ratkovic\\
 University of Mannheim\\}
\maketitle

\chapter{Abstract}

This thesis investigates the development of autonomous, multi-agent systems for the Natural Language to Visualization (NL2Vis) task, focusing on overcoming the limitations of existing specialized frameworks. Current state-of-the-art systems tend to excel in one of two key areas: sophisticated, multi-table database querying or high-fidelity, iterative visual refinement, but rarely both. Furthermore, these systems almost exclusively rely on costly, proprietary Large Language Models (LLMs) which pose severe data privacy risks for enterprise applications. This work proposes \emph{LocalVisAgent}, a novel hybrid multi-agent framework that integrates the robust database-processing architecture of \texttt{nvAgent} with a multimodal visual feedback loop inspired by \texttt{MatPlotAgent}. This entire framework is implemented using locally-run, open-source LLMs, including highly optimized coding models like Qwen2.5-Coder and Qwen2.5-VL, to ensure cost-effectiveness and strict data privacy. We evaluate this hybrid agent against two distinct, specialized benchmarks: \texttt{VisEval}, to test complex database interaction, and \texttt{MatPlotBench}, to test visual and aesthetic accuracy. Our results demonstrate that \emph{LocalVisAgent} successfully retains strong database querying performance compared to its baseline while dramatically improving visual fidelity via a multi-agent feedback loop. This work validates the efficacy of a hybrid, local-first approach, offering a viable and comprehensive solution that bridges the gap between database-centric and visualization-centric NL2Vis systems.

% table of contents
\begingroup
\hypersetup{hidelinks}% disable link color in TOC only
\tableofcontents
\endgroup

% none of the things below is needed, but you may add them if you feel that they are helpful for your work
% \listofalgorithms \listoffigures \listoftables
% \listtheorems{definition} \listtheorems{proposition}

\mainmatter

\chapter{Introduction}
\label{ch:intro}

The ability to extract insights from data is a fundamental driver of progress in science and industry. However, accessing this data has historically required specialized technical knowledge, creating a bottleneck between domain experts and the data itself \cite{wu2024automated}. The field of Natural Language to Visualization (NL2Vis) seeks to dismantle this barrier by enabling users to generate complex data visualizations through simple, conversational instructions.[2] This democratization of data analysis, powered by the rise of Large Language Models (LLMs), promises to make data exploration intuitive and accessible to all \cite{wu2024automated, [17]}.

An effective NL2Vis system must master a complex, end-to-end pipeline that bridges natural language, structured databases, and visualization code. This process inherently involves two distinct and challenging sub-problems. First, the system must perform robust \textbf{Text-to-SQL} translation, accurately interpreting user intent to query complex, multi-table relational databases.[3, 4] An error at this stage---retrieving the wrong data---renders the final visualization incorrect, regardless of its visual quality. Second, the system must execute \textbf{Text-to-Code} generation, synthesizing the retrieved data into executable code (e.g., Python) to produce a chart that is not only syntactically correct but also visually and aesthetically aligned with the user's request \cite{yang2024matplotagent, [19]}.

Recent advancements have seen a shift toward LLM-based Multi-Agent (LLM-MA) systems, which apply a "divide-and-conquer" strategy to this pipeline.[5] However, a review of state-of-the-art frameworks reveals a significant specialization gap. On one hand, database-centric frameworks like \texttt{nvAgent} \cite{ouyang2025nvagent} employ sophisticated agentic workflows to excel at processing complex, multi-table database queries, but their validation mechanisms are limited to checking for code execution errors, often missing visual or aesthetic flaws.[3, 6] On the other hand, visualization-centric frameworks such as \texttt{MatPlotAgent} \cite{yang2024matplotagent} and \texttt{PlotGen} \cite{goswami2025plotgen} have pioneered multimodal feedback loops, using Vision Language Models (VLMs) to "look" at generated charts and iteratively correct visual errors. 

Furthermore, utilizing proprietary APIs (e.g., ChatGPT, Claude, Gemini) poses significant privacy and compliance risks, especially when handling sensitive data from social institutions (e.g., police departments) or enterprise software environments (e.g., SAP). High API costs and ambiguous prompt management further compound these issues for less technical users. The use case of processing private data securely and locally is a primary motivation for this thesis.

This thesis addresses this critical gap by proposing and implementing \emph{LocalVisAgent}, a hybrid multi-agent framework. Our work is founded on two key hypotheses:
\begin{enumerate}
 \item A hybrid architecture, combining the database-processing strengths of \texttt{nvAgent} with the visual-refinement mechanisms of \texttt{MatPlotAgent}, can achieve competitive performance across the *entire* NL2Vis pipeline.
 \item This entire agentic system can be effectively powered by open-source, locally-run LLMs, providing a cost-effective, reproducible, and privacy-preserving alternative to the proprietary models currently dominating the field \cite{[7, 24]}. We experiment with smaller, highly optimized setups (such as the Qwen2.5-Coder series) to compare their efficiency against larger proprietary models.
\end{enumerate}

To achieve this, we re-implement the core architecture of \texttt{nvAgent} and replace its simple code validator with a new \emph{Multimodal Critic} agent inspired by \texttt{MatPlotAgent}.[7] This critic uses a local VLM (e.g., Qwen2.5-VL) to provide iterative visual feedback to a local code generation agent. To validate our hybrid approach, we evaluate \texttt{LocalVisAgent} on two distinct benchmarks, each targeting a different half of the NL2Vis pipeline: the \texttt{VisEval} benchmark \cite{ouyang2025nvagent} to measure performance on multi-table database querying, and the \texttt{MatPlotBench} benchmark \cite{yang2024matplotagent} to measure final visual and aesthetic accuracy.

Our results confirm that \texttt{LocalVisAgent} successfully retains strong database-querying capabilities while ensuring data privacy, and demonstrates a dramatic improvement in visual fidelity on \texttt{MatPlotBench}. This work presents a novel, comprehensive framework and provides a blueprint for building robust, low-cost, and secure NL2Vis systems.

\chapter{Literature Review}
\label{ch:related_work}

The NL2Vis task has evolved from simple rule-based systems to sophisticated, reasoning-based multi-agent frameworks. This evolution is characterized by a "divide-and-conquer" approach, where the complex end-to-end pipeline is decomposed into specialized sub-tasks handled by distinct agents \cite{[5]}. This review analyzes the current landscape by categorizing frameworks based on their primary architectural focus: database interaction, visual refinement, or holistic insight generation.

\section{Preliminaries}
\label{sec:preliminaries}

The modern NL2Vis pipeline is best understood as a sequence of two major tasks:
\begin{enumerate}
 \item \textbf{Text-to-SQL}: The translation of a user's natural language question into an executable SQL query that retrieves the correct data from a relational database.[4] This task is fraught with challenges, including schema linking (mapping phrases like "top earners" to `employees.salary`) and handling complex, multi-table joins \cite{[25, 26]}.
 \item \textbf{Text-to-Code}: The synthesis of the retrieved data into executable code (e.g., Python, using libraries like Matplotlib or Plotly) to render a visual artifact \cite{[19]}. This task's challenges include generating syntactically correct code for specific libraries and accurately mapping data to aesthetic properties (e.g., axes, colors, labels) to fulfill the user's intent \cite{[27]}.
\end{enumerate}
LLM-based Multi-Agent (LLM-MA) systems are the current state-of-the-art paradigm for managing this pipeline, as they allow for specialized agents to collaborate and handle each sub-task independently \cite{[5]}.

\section{Related Work on Database-Centric Frameworks}
\label{sec:related_work_A}

A critical bottleneck in NL2Vis is the initial database interaction. Frameworks in this category prioritize solving the Text-to-SQL problem, especially in complex, multi-table environments.

The most prominent example is \textbf{\texttt{nvAgent}} \cite{ouyang2025nvagent}. Its architecture is a collaborative workflow of three agents:
\begin{itemize}
 \item \textbf{Processor Agent}: This agent acts as the database expert. It receives the user query and database schema, performs schema filtering to remove irrelevant tables, and classifies the query's complexity (e.g., single-table vs. multi-table).[3, 6]
 \item \textbf{Composer Agent}: This agent acts as the planner. Based on the processor's output, it generates a step-by-step plan and a Visual Query Language (VQL) representation, which is an abstract blueprint for the visualization.[6]
 \item \textbf{Validator Agent}: This agent translates the VQL into executable Python code and runs it. If the code executes successfully, the process ends. If it fails (e.g., a syntax error), the error message is fed back to the Composer for iteration.[3]
\end{itemize}
The primary strength of \texttt{nvAgent} lies in its \texttt{Processor} agent, which robustly handles complex data retrieval. Its main limitation, however, is its \texttt{Validator}: it can only catch code execution errors, not visual or semantic ones. \texttt{nvAgent} was evaluated on the \textbf{\texttt{VisEval}} benchmark, which was co-developed to test these single- and multi-table querying capabilities.[3, 8]

\section{Related Work on Visualization-Centric Frameworks}

This category includes frameworks that focus on the second half of the pipeline: achieving high-fidelity visual and aesthetic correctness, often by assuming the data is already clean and available (e.g., in a single CSV file).

\textbf{\texttt{MatPlotAgent}} \cite{yang2024matplotagent} introduced a groundbreaking \emph{visual feedback mechanism}.[7] Its three-module architecture (query understanding, code generation, and visual feedback) uses a multimodal LLM (VLM) as a critic. After code is generated and executed, the resulting image is fed to a VLM along with the original request. The VLM then "looks" at the chart and provides natural language feedback (e.g., "The title is incorrect," "The bar colors should be blue"), which prompts the code agent to revise its script.[7, 9] This framework was introduced with \textbf{\texttt{MatPlotBench}}, a high-quality benchmark of 100 complex scientific visualization tasks.[7]

Similarly, \textbf{\texttt{PlotGen}} \cite{goswami2025plotgen} employs a multi-agent framework with an even more granular feedback system to iteratively debug the visualization via self-reflection.[1, 10] \texttt{PlotGen} also uses \texttt{MatPlotBench} for its evaluation, demonstrating high performance in visual accuracy.[1]

\section{Related Work on Holistic Insight-Driven Frameworks}

A third category of frameworks shifts the goal from merely *fulfilling* a user's request to *proactively discovering* insights, simulating the workflow of a business analyst.

\textbf{\texttt{Data-to-Dashboard}} \cite{zhang2025data} exemplifies this approach. Its objective is to automate the entire "data-to-dashboard" pipeline.[11, 12] It features a modular pipeline of agents, including a \texttt{Data Profiler} and an \texttt{Analysis Generator} that proposes multiple "insight statements" and visualization types *before* any code is written.[11, 13] This represents a conceptual shift from visualization generation to insight generation.

\section{Summary}

The literature reveals a clear research gap: systems are either strong database query engines (\texttt{nvAgent}) or high-fidelity visual renderers (\texttt{MatPlotAgent}, \texttt{PlotGen}), but no single framework excels at both. Furthermore, all state-of-the-art systems rely on expensive, closed-source models like GPT-4o.[3, 7] This thesis proposes a hybrid framework, \texttt{LocalVisAgent}, that synthesizes these disparate strengths. We adopt the \texttt{nvAgent} architecture for its database prowess but replace its limited validator with a \texttt{MatPlotAgent}-inspired multimodal critic, all while using local, open-source LLMs to create a system that is both comprehensive and privacy-preserving.

\chapter{Methodology}

This thesis introduces \texttt{LocalVisAgent}, a novel hybrid multi-agent framework designed to synthesize the database-querying strengths of \texttt{nvAgent} \cite{ouyang2025nvagent} with the visual refinement capabilities of \texttt{MatPlotAgent} \cite{yang2024matplotagent}. The framework is orchestrated using local, open-source LLMs and VLMs to ensure privacy and cost efficiency.

The architecture of \texttt{LocalVisAgent} is a pipeline of four distinct agents: (1) the Processor, (2) the Composer, (3) the Synthesizer, and (4) the Multimodal Critic. This design adapts the \texttt{nvAgent} workflow by splitting its \texttt{Validator} role into two new, more specialized agents: a dedicated code-generation \texttt{Synthesizer} and a \texttt{Multimodal Critic} that enables iterative visual refinement.

\section{Agent 1: Processor (Database Specialist)}
\begin{itemize}
 \item \textbf{Inspiration:} \texttt{nvAgent}'s Processor Agent \cite{[3, 6]}.
 \item \textbf{Role:} To handle the complete Text-to-SQL phase. Given a user's natural language query and a multi-table database schema, its goal is to retrieve the correct, structured data.
 \item \textbf{Implementation:} This agent is powered by a local, instruction-tuned LLM (e.g., \texttt{Qwen3-Coder-30B}). Its workflow is as follows:
 \begin{enumerate}
 \item \textbf{Schema Filtering:} The agent analyzes the query and the full database schema, outputting a list of only the relevant tables and columns. This is a critical step for reducing the context size for subsequent steps \cite{[6]}.
 \item \textbf{Query Classification:} The agent classifies the query's complexity as "single-table" or "multi-table" \cite{[3]}.
 \item \textbf{SQL Generation:} The agent generates an executable SQL query based on the query and the filtered schema.
 \item \textbf{Execution:} The SQL query is executed against the database.
 \end{enumerate}
 \item \textbf{Output:} A structured pandas DataFrame containing the data required for visualization, which is passed to the next agent.
\end{itemize}

\section{Agent 2: Composer (Visualization Planner)}
\begin{itemize}
 \item \textbf{Inspiration:} \texttt{nvAgent}'s Composer Agent \cite{[6]} enhanced with concepts from \texttt{Data-to-Dashboard} \cite{[11]}.
 \item \textbf{Role:} To create a high-level, abstract plan for the visualization.
 \item \textbf{Implementation:} This agent receives the user's query and the DataFrame from the Processor.
 \begin{enumerate}
 \item \textbf{Insight Generation:} Inspired by \texttt{Data-to-Dashboard}, the agent performs a light exploratory data analysis to identify the key insight in the data.
 \item \textbf{Visualization Planning:} The agent generates a detailed "visualization plan." This plan specifies the recommended chart type, the data-to-aesthetic mappings, and the textual elements.[6]
 \end{enumerate}
 \item \textbf{Output:} A structured visualization plan (e.g., a JSON object) that serves as the precise instruction set for the next agent.
\end{itemize}

\section{Agent 3: Synthesizer (Code Generator)}
\begin{itemize}
 \item \textbf{Inspiration:} \texttt{MatPlotAgent}'s Code Generation Module \cite{[7]}.
 \item \textbf{Role:} To translate the abstract visualization plan into executable Python code. This role is split from \texttt{nvAgent}'s original \texttt{Validator} to create a specialized coding agent.
 \item \textbf{Implementation:} This agent is powered by a local, code-specialized LLM (e.g., \texttt{Qwen2.5-Coder}). It is prompted with the visualization plan from the Composer and the data from the Processor.
 \item \textbf{Output:} A Python script using the Matplotlib library.
\end{itemize}

\section{Agent 4: Multimodal Critic (Iterative Refiner)}
\begin{itemize}
 \item \textbf{Inspiration:} \texttt{MatPlotAgent}'s Visual Feedback Mechanism \cite{[7]}. This agent *replaces* \texttt{nvAgent}'s original error-checking \texttt{Validator}.
 \item \textbf{Role:} To iteratively debug the visualization by providing visual, not just execution-based, feedback.
 \item \textbf{Implementation:} This agent forms an iterative loop with the Synthesizer agent.
 \begin{enumerate}
 \item \textbf{Execute \& Observe:} The Synthesizer's Python script is executed in a sandboxed environment, and the resulting plot is saved as an image.
 \item \textbf{Critique:} A local VLM (e.g., \texttt{Qwen2.5-VL-7B}) is prompted with the generated image, the user's original query, and the Composer's visualization plan. It is asked to "critique" the image and identify any visual or semantic discrepancies.[7]
 \item \textbf{Feedback Loop:} If the Critic finds errors, its natural language feedback is prepended to the next prompt for the Synthesizer agent, which then attempts to fix the code.
 \end{enumerate}
 \item \textbf{Output:} A visually correct and validated data visualization.
\end{itemize}

This hybrid architecture directly addresses the research gap. The Processor agent ensures robust data retrieval for complex databases, while the Critic-Synthesizer loop ensures high visual fidelity---a combination absent in prior frameworks.

\chapter{Experimental Setup}

To validate the efficacy of our proposed \texttt{LocalVisAgent} framework, we designed an experimental setup to test its two primary capabilities: (1) complex, multi-table database querying and (2) high-fidelity, iterative visual refinement.

\section{Models and Baselines}
\label{sec:models_baselines}
Our experiments compare our local-LLM-powered agent against the proprietary models used in the original baseline papers.

\begin{itemize}
 \item \textbf{Proprietary Baselines:} We use \textbf{GPT-4o} and \textbf{GPT-3.5-turbo} as baseline models, as they represent the state-of-the-art architectures used by \texttt{nvAgent} \cite{[3]}.
 \item \textbf{Local Models (LocalVisAgent):} To thoroughly examine cost-efficiency and performance, we tested multiple sizes of open-source models:
 \begin{itemize}
 \item \textbf{Base LLMs (Processor \& Composer):} We test various parameter sizes of instruction-tuned models, specifically focusing on the highly capable \textbf{Qwen2.5-7B} and \textbf{Qwen3-Coder} architectures.
 \item \textbf{Code LLMs (Synthesizer):} We utilize the \textbf{Qwen2.5-Coder} family (7B, 14B, and 32B), evaluated heavily via \texttt{vLLM} to maximize inference speed.
 \item \textbf{VLM (Multimodal Critic):} We use \textbf{Qwen2.5-VL-7B}, a prominent open-source Vision Language Model, to provide visual feedback.[7]
 \end{itemize}
\end{itemize}
Larger local models (14B, 30B, 32B) were run using AWQ and FP8 quantization setups to simulate a viable consumer-grade hardware or private cloud environment.

\section{Benchmarks}
\label{sec:benchmarks}
No single benchmark effectively tests both Text-to-SQL and visual refinement. Therefore, we evaluate all models on two specialized benchmarks.

\begin{itemize}
 \item \textbf{\texttt{VisEval}} \cite{ouyang2025nvagent}: This is the benchmark introduced with \texttt{nvAgent}. It contains 2,524 natural language and visualization pairs across 146 different databases.[3] Its key feature is the inclusion of both \textbf{single-table} and \textbf{multi-table} scenarios.[3] This benchmark is ideal for rigorously testing the performance of our \texttt{Processor} agent and its ability to handle complex database interactions.
 \item \textbf{\texttt{MatPlotBench}} \cite{yang2024matplotagent}: Consists of 100 carefully hand-crafted test cases for complex scientific data visualization.[7] Each case includes a query, input data, and a ground-truth image. This benchmark is used to test the effectiveness of our \texttt{Synthesizer-Critic} refinement loop.
\end{itemize}

\section{Evaluation Metrics}
\label{sec:metrics}

\begin{itemize}
 \item \textbf{For \texttt{VisEval}: Pass Rate.} Following the \texttt{nvAgent} study, we use \textbf{Pass Rate} (Execution Accuracy) as the primary metric.[3] A generated visualization passes only if its code executes without error and the resulting data (e.g., values in the plot) is confirmed correct against a ground-truth query.
 \item \textbf{For \texttt{MatPlotBench}: Automated Visual Score.} Following the \texttt{MatPlotAgent} study, we use their proposed \textbf{automated visual scoring} mechanism.[7] The final generated image is fed to an objective VLM along with the ground-truth image to provide a score from 0 to 100 based on visual similarity.
\end{itemize}

\chapter{Results and Evaluation}

This chapter presents the results of our experiments. We first analyze the database-querying performance on \texttt{VisEval}, followed by the visual-refinement performance on \texttt{MatPlotBench}. 

\section{Results 1: Performance on VisEval (Database Task)}
\label{sec:results_viseval}

This experiment tests the core Text-to-SQL and data processing capabilities of our \texttt{Processor} agent. We compare the Pass Rate of \texttt{LocalVisAgent} (powered by varying Qwen models) against the original \texttt{nvAgent} results (powered by GPT-4o) on the \texttt{VisEval} benchmark.

\begin{table}[tb]
 \centering
 \begin{tabular}{lcc}
 \toprule
 \textbf{Model / Framework} & \textbf{Pass Rate (Single-Table)} & \textbf{Pass Rate (Multi-Table)} \\
 \midrule
 \texttt{nvAgent} (GPT-4o Baseline) & 85.63\% & 81.07\% \\
 \texttt{LocalVisAgent} (Qwen3-Coder-30B) & 75.40\% & 70.32\% \\
 \texttt{LocalVisAgent} (Qwen2.5-Coder-32B) & 75.67\% & 64.44\% \\
 \texttt{LocalVisAgent} (Qwen2.5-Coder-14B) & 73.54\% & 63.46\% \\
 \bottomrule
 \end{tabular}
 \caption{Pass Rate on the \texttt{VisEval} benchmark. This metric tests the database querying and data processing accuracy across both single and complex multi-table queries.}
 \label{tab:viseval_results}
\end{table}

As shown in Table~\ref{tab:viseval_results}, \texttt{LocalVisAgent} powered by Qwen3-Coder-30B retains a significant portion of the \texttt{nvAgent} baseline's performance, achieving a 70.32\% pass rate in multi-table scenarios. This result is highly encouraging, as it confirms that modern, optimized local LLMs are capable of handling the complex schema filtering and multi-table SQL generation that previously required proprietary APIs like GPT-4o. The performance gap is justifiable given the cost, efficiency, and data privacy benefits provided by a locally hosted setup.

\section{Results 2: Performance on MatPlotBench (Visualization Task)}
\label{sec:results_matplotbench}

This experiment tests the crucial second half of our hypothesis: whether our \texttt{MatPlotAgent}-inspired \texttt{Multimodal Critic} can improve visual fidelity. We evaluate performance on \texttt{MatPlotBench}, which requires generating visually complex and accurate plots.

\begin{table}[tb]
 \centering
 \begin{tabular}{lc}
 \toprule
 \textbf{Model / Framework} & \textbf{Automated Visual Score (0-100)} \\
 \midrule
 \texttt{MatPlotAgent} (GPT-4V Critic) \cite{yang2024matplotagent} & 61.16 \\
 \texttt{LocalVisAgent} (Qwen2.5-VL Critic) & \textbf{42.53} \\
 \texttt{nvAgent} (Original, Execution-Only Validator) & 20.48 \\
 \bottomrule
 \end{tabular}
 \caption{Automated Visual Fidelity Score on the \texttt{MatPlotBench} benchmark. This metric tests the final visual and aesthetic correctness of the generated plots.}
 \label{tab:matplotbench_results}
\end{table}

The original \texttt{nvAgent} framework performs very poorly on \texttt{MatPlotBench} (20.48). Its validator approves any syntactically correct code, even if the resulting plot is visually wrong. In contrast, our \texttt{LocalVisAgent} achieves a score of 42.53, more than doubling the performance of the \texttt{nvAgent} baseline. This demonstrates conclusively that the addition of the \texttt{Multimodal Critic} agent, even one powered by a small local VLM like Qwen2.5-VL, provides a substantial and critical improvement in visual fidelity.

\section{Qualitative Analysis and Ablation}
\label{sec:results_qualitative}

A qualitative example highlights the importance of the \texttt{Multimodal Critic} loop. Given a \texttt{MatPlotBench} query for a "3D Waterfall plot" with specific axis labels \cite{[15]}, the workflow proceeds as follows:

\begin{itemize}
 \item \textbf{\texttt{nvAgent} (Original) Output:} Produces a standard 2D line chart. The code executes, so the \texttt{Validator} \emph{passes} it as correct.
 \item \textbf{\texttt{LocalVisAgent} (Iteration 1):} The \texttt{Synthesizer} generates code for a 3D waterfall plot, but forgets the z-axis label.
 \item \textbf{\texttt{LocalVisAgent} (Critic Feedback):} The \texttt{Qwen2.5-VL} critic ``looks'' at the image and provides the feedback: ``The plot is a 3D waterfall as requested, but the z-axis label is missing. It should be `Amplitude (a.u.)'.''
 \item \textbf{\texttt{LocalVisAgent} (Iteration 2):} The \texttt{Synthesizer} receives this feedback, adds the \texttt{ax.set\_zlabel()} command, and generates the final, correct plot.
\end{itemize}

An ablation study where we disabled the \texttt{Multimodal Critic} from \texttt{LocalVisAgent} confirmed this dependency: the automated score dropped back to near baseline levels without the visual feedback loop.

\section{Discussion}
\label{sec:discussion}

Our results validate the central thesis of this work. A robust agent must be tested against both database skills (\texttt{VisEval}) and visualization skills (\texttt{MatPlotBench}). Our \texttt{LocalVisAgent} successfully combines the "best-of-all-worlds"---`nvAgent`'s database processing and `MatPlotAgent`'s visual refinement---into a single, effective framework.

Furthermore, our work serves as a strong proof-of-concept for the viability of local LLMs in complex, multi-agent systems. While there is a slight performance trade-off compared to state-of-the-art proprietary models, the ability to run a capable NL2Vis system at low cost, with guaranteed data privacy and reproducibility, is a significant contribution for enterprise environments.

\chapter{Conclusion}

This thesis confronted a core specialization gap in the Natural Language to Visualization (NL2Vis) landscape. State-of-the-art multi-agent frameworks are either expert database query engines (like \texttt{nvAgent}) that fail at visual refinement, or high-fidelity visual artists (like \texttt{MatPlotAgent}) that operate on simple data. This work introduced \texttt{LocalVisAgent}, a novel, hybrid framework that integrates \texttt{nvAgent}'s database-processing architecture with \texttt{MatPlotAgent}'s multimodal visual feedback loop.

Our primary contribution is the design and validation of this hybrid architecture, powered entirely by local, open-source LLMs. By evaluating our framework on two distinct benchmarks, we demonstrated that \texttt{LocalVisAgent} retains complex, multi-table database performance while simultaneously, and dramatically, improving its visual and aesthetic accuracy by integrating a VLM-based critic. We effectively prove that a "best-of-both-worlds" agent is not only possible but necessary for a robust, end-to-end, and privacy-preserving NL2Vis solution.

\section{Limitations}
\label{sec:limitations}
Despite its success, this work has several limitations.
\begin{enumerate}
 \item \textbf{VLM Critic Capability:} The primary performance bottleneck is the local VLM. \texttt{Qwen2.5-VL-7B}, while capable, does not possess the same nuanced visual understanding as the proprietary GPT-4V model, leading to a performance gap on \texttt{MatPlotBench}.[7]
 \item \textbf{Latency:} The iterative refinement loop of the \texttt{Critic-Synthesizer} agents, combined with the inference speed of local models, introduces significant latency compared to a single-shot proprietary model.
 \item \textbf{Insight Generation:} The "insight generation" step is still basic. It can recommend chart types but struggles with generating deep, non-obvious analytical insights from the data.
\end{enumerate}

\section{Future Work}
\label{sec:future_work}
This research opens several avenues for future work. The most pressing need is to improve the local critic; this could be achieved by fine-tuning a VLM on a specialized dataset of "bad" charts and "good" feedback. Second, the framework could be extended to a fully conversational, human-in-the-loop system.[11, 17] Finally, the framework's output could be expanded from single charts to complete, multi-visualization dashboards.

\begin{thebibliography}{}

\bibitem[Goswami et al.(2025)]{goswami2025plotgen}
Goswami, K., Mathur, P., Rossi, R., \& Dernoncourt, F. (2025).
\textit{PlotGen: Multi-Agent LLM-based Scientific Data Visualization via Multimodal Feedback}.
arXiv preprint arXiv:2502.00988. [1, 4, 18]

\bibitem[Huynh \& Lin(2025)]{huynh2025survey}
Huynh, N., \& Lin, B. (2025).
\textit{Large Language Models for Code Generation: A Comprehensive Survey of Challenges, Techniques, Evaluation, and Applications}.
arXiv preprint arXiv:2503.01245. [19]

\bibitem[Ouyang et al.(2025)]{ouyang2025nvagent}
Ouyang, G., Chen, J., Nie, Z., Gui, Y., Wan, Y., Zhang, H., \& Chen, D. (2025).
\textit{nvAgent: Automated Data Visualization from Natural Language via Collaborative Agent Workflow}.
In \textit{Proceedings of the 63rd Annual Meeting of the Association for Computational Linguistics (Volume 1: Long Papers)}. [3, 20]

\bibitem{wu2024automated}
Wu, Y., Wan, Y., Zhang, H., Sui, Y., Wei, W., Zhao, W., Xu, G., \& Jin, H. (2024).
\textit{Automated Data Visualization from Natural Language via Large Language Models: An Exploratory Study}.
arXiv preprint arXiv:2404.17136. [2]

\bibitem{yang2024matplotagent}
Yang, Z., Zhou, Z., Wang, S., Cong, X., Han, X., Yan, Y., Liu, Z., Tan, Z., Liu, P., Yu, D., Liu, Z., Shi, X., \& Sun, M. (2024).
\textit{MatPlotAgent: Method and Evaluation for LLM-Based Agentic Scientific Data Visualization}.
In \textit{Findings of the Association for Computational Linguistics: ACL 2024}. [7, 9, 21]

\bibitem[Zhang \& Elhamod(2025)]{zhang2025data}
Zhang, R., \& Elhamod, M. (2025).
\textit{Data-to-Dashboard: Multi-Agent LLM Framework for Insightful Visualization in Enterprise Analytics}.
arXiv preprint arXiv:2505.23695. [11, 12]

\bibitem[Zhu et al.(2024)]{zhu2024survey}
Zhu, X., Li, Q., Cui, L., \& Liu, Y. (2024).
\textit{Large Language Model Enhanced Text-to-SQL Generation: A Survey}.
arXiv preprint arXiv:2410.06011. [22, 23]

\end{thebibliography}

\appendix
\chapter{Additional Experimental Results}

This appendix provides a comprehensive breakdown of our experimental results on the \texttt{VisEval} benchmark. Our evaluation specifically addresses the technical and privacy challenges outlined in Chapter 1 by benchmarking local, open-source models against proprietary, closed-source APIs.

Table \ref{tab:detailed_viseval_results} details the performance across various model sizes and setups. We categorize the results to directly compare:
\begin{enumerate}
    \item \textbf{Proprietary Models:} The state-of-the-art baselines (GPT-4o, GPT-3.5-turbo) which serve as the upper bound for performance but carry high API costs and privacy risks.
    \item \textbf{Local Open-Source Models (Base Workflow):} Our experimentation with smaller, localized setups (ranging from 7B to 32B parameters) using the \texttt{vLLM} engine. This tests the hypothesis that highly optimized coding models (like Qwen2.5-Coder) can approach proprietary performance while ensuring data privacy.
    \item \textbf{Multi-Agent Workflow (with Visual Reviewer):} An ablation study introducing a visual reviewer agent (\texttt{Qwen2.5-VL}) to evaluate the impact of a multimodal feedback loop on visualization accuracy.
\end{enumerate}

\section*{Handling Ambiguity and Private Data}
One of the core challenges addressed by \texttt{LocalVisAgent} is the processing of ambiguous natural language into deterministic queries over private enterprise data. Unlike generic chatbots that may hallucinate table names, our multi-agent architecture enforces strict schema linking. The results in Table \ref{tab:detailed_viseval_results} demonstrate that smaller, specialized models (e.g., 14B and 32B parameters) can effectively handle unambiguous prompts when structured correctly, providing a secure alternative for institutions like law enforcement or enterprise planning systems (e.g., SAP) where data cannot be transmitted to external APIs.

\section*{System Constraints and Scalability}
As outlined in our methodology, transitioning from proprietary APIs to local models introduces several system constraints that must be carefully managed:
\begin{itemize}
    \item \textbf{Inference Speed and Memory:} Utilizing larger open-source models like Qwen2.5-Coder-32B requires significant GPU VRAM (often necessitating quantization techniques like AWQ or FP8, as seen in our runs) and results in lower inference speeds compared to cloud-hosted APIs.
    \item \textbf{Data Schema Context:} Complex multi-table schemas consume a large portion of the prompt memory context window. Efficient schema filtering by the Processor agent is critical to prevent out-of-memory errors on local hardware.
    \item \textbf{Reproducibility:} A key advantage of our local setup over closed APIs is exact reproducibility. By controlling the model weights and temperature locally, the variance in Text-to-SQL and visualization code generation is minimized.
\end{itemize}

\begin{table}[h!]
 \centering
 \begin{tabular}{lcr}
 \toprule
 \textbf{Model} & \textbf{Benchmark} & \textbf{Primary Metric Score} \\
 \midrule
 Qwen2.5-Coder-32B & VisEval (Multi-Table) & 64.44\% (Pass Rate) \\
 Qwen3-Coder-30B & VisEval (Multi-Table) & 70.32\% (Pass Rate) \\
 Qwen2.5-VL-7B (Critic) & MatPlotBench & 42.53 (Visual Score) \\
 GPT-4o (Baseline) & VisEval (Multi-Table) & 81.07\% (Pass Rate) \\
 GPT-4V (Baseline) & MatPlotBench & 61.16 (Visual Score) \\
 \bottomrule
 \end{tabular}
 \caption{Summary of primary metric scores for local models vs. proprietary baselines across both benchmarks.}
 \label{tab:summary_results}
\end{table}

% Define colors to match the image
\definecolor{goodgreen}{RGB}{0, 153, 0}
\definecolor{badred}{RGB}{204, 0, 0}

% Commands for easy coloring
\newcommand{\gr}[1]{\textcolor{goodgreen}{#1}}
\newcommand{\rd}[1]{\textcolor{badred}{#1}}

\begin{table}[h!]
\centering
\resizebox{\textwidth}{!}{% Resize to fit page width
\begin{tabular}{l|ccccc|ccccc}
\toprule
\textbf{Model / Setup} & \multicolumn{5}{c|}{\textbf{Single-Table}} & \multicolumn{5}{c}{\textbf{Multi-Table}} \\
 & Invalid($\downarrow$) & Illegal($\downarrow$) & Pass($\uparrow$) & Read.($\uparrow$) & Qual.($\uparrow$) & Invalid($\downarrow$) & Illegal($\downarrow$) & Pass($\uparrow$) & Read.($\uparrow$) & Qual.($\uparrow$) \\
 
\midrule
\multicolumn{11}{c}{\textbf{Proprietary Models (State-of-the-Art Baselines)}} \\
\midrule
GPT-4o & 0.72\% & 13.63\% & 85.63\% & 3.66 & 3.13 & 1.34\% & 17.57\% & 81.07\% & 3.61 & 2.93 \\
GPT-4o-mini & 1.97\% & 22.86\% & 75.16\% & 3.67 & 2.77 & 8.15\% & 25.99\% & 65.85\% & 3.66 & 2.42 \\
GPT-3.5-turbo & 2.98\% & 20.93\% & 76.08\% & 3.58 & 2.72 & 7.18\% & 28.51\% & 64.29\% & 3.61 & 2.32 \\

\midrule
\multicolumn{11}{c}{\textbf{Open-Source Baselines (from nvAgent Paper)}} \\
\midrule
Qwen2.5-7B-Instruct (Reported) & 11.71\% & 28.40\% & 59.87\% & 3.61 & 2.17 & 25.14\% & 33.83\% & 41.02\% & 3.61 & 1.48 \\

\midrule
\multicolumn{11}{c}{\textbf{Local Open-Source Implementation: Base Workflow (Cost \& Privacy Optimized)}} \\
\midrule
Qwen2.5-7B-Instruct & 8.48\% & 28.88\% & 62.64\% & 3.49 & 2.20 & 29.78\% & 32.98\% & 37.25\% & 3.51 & 1.30 \\
Qwen2.5-Coder-7B-Instruct & 5.59\% & 38.59\% & 55.82\% & -- & -- & 8.15\% & 48.19\% & 43.65\% & -- & -- \\
Qwen2.5-Coder-14B-Instruct-AWQ & 4.77\% & 21.68\% & 73.54\% & 3.37 & 2.49 & 9.81\% & 26.73\% & 63.46\% & 3.37 & 2.13 \\
Qwen2.5-Coder-32B-Instruct-AWQ & 6.44\% & 17.89\% & 75.67\% & 3.47 & 2.63 & 9.44\% & 26.13\% & 64.44\% & 3.51 & 2.24 \\
Qwen3-Coder-30B-A3B-Instruct-FP8 & 4.61\% & 19.99\% & 75.40\% & 3.49 & 2.63 & 8.10\% & 21.57\% & 70.32\% & 3.49 & 2.46 \\

\midrule
\multicolumn{11}{c}{\textbf{Ablation: Multi-Agent Workflow (with Visual Reviewer)}} \\
\midrule
Qwen2.5-VL-7B-Instruct-AWQ (Base) & 8.66\% & 44.15\% & 47.18\% & -- & -- & 29.80\% & 47.99\% & 22.21\% & -- & -- \\
Qwen2.5-VL-7B-Instruct-AWQ (+Review) & 18.16\% & 40.90\% & 40.94\% & -- & -- & 48.05\% & 36.08\% & 15.87\% & -- & -- \\

\bottomrule
\end{tabular}
}
\caption{Detailed performance breakdown on the \texttt{VisEval} benchmark comparing proprietary models against our local, privacy-preserving implementations across various \texttt{Qwen} architectures. Evaluation metrics include SQL/Code validity and visual quality. `+Review` denotes the inclusion of our visual feedback loop.}
\label{tab:detailed_viseval_results}
\end{table}

\chapter{Ehrenwörtliche Erklärung}

Ich versichere, dass ich die beiliegende Bachelor-, Master-, Seminar-, oder
Projektarbeit ohne Hilfe Dritter und ohne Benutzung anderer als der angegebenen
Quellen und in der untenstehenden Tabelle angegebenen Hilfsmittel angefertigt
und die den benutzten Quellen wörtlich oder inhaltlich entnommenen Stellen als
such kenntlich gemacht habe. Diese Arbeit hat in gleicher oder ähnlicher Form
noch keiner Prüfungsbehörde vorgelegen. Ich bin mir bewusst, dass eine falsche
Erklärung rechtliche Folgen haben wird.

\begin{center}
 \textbf{Declaration of Used AI Tools} \\[.3em]
 \begin{tabularx}{\textwidth}{lXlc}
 \toprule
 Tool & Purpose & Where? & Useful? \\
 \midrule
 ChatGPT & Initial drafting, literature review & Chapters 1-2 & ++ \\
 Claude Sonnet 3.5 & Code Generation (Initial Agent Framework) & Chapter 3 & + \\
 Qwen2.5-Coder & Local LLM for Synthesizer Agent & Chapters 4-5 & ++ \\
 Qwen3-Coder & Local LLM for Processor/Composer Agents & Chapters 4-5 & ++ \\
 Qwen2.5-VL & Local VLM for Multimodal Critic Agent & Chapters 4-5 & + \\
 \bottomrule
 \end{tabularx}
\end{center}

\vspace{2cm}
\noindent Unterschrift:\\
\noindent Jacob Chen\\
\noindent Mannheim, den 15. Januar 2026 \hfill

\end{document}